% ============================================================
% CAPITOLO 5 — Progetto 2: Applicazione CLI Completa
% ============================================================
\chapter{Progetto 2: Un'Applicazione CLI Completa}\index{CLI}

\section*{Cosa Costruirai}
Un \textbf{Task Manager da riga di comando} --- un'applicazione Python professionale che:
\begin{itemize}
  \item Gestisce task con operazioni CRUD (crea, leggi, aggiorna, elimina)
  \item Persiste i dati su file JSON
  \item Supporta filtri e ricerca
  \item Ha un sistema di priorità e scadenze
  \item Include test automatici con copertura completa
\end{itemize}

\textbf{Tempo stimato}: 30--45 minuti

% ---------------------------------------------------------
\section{Il Contesto Professionale (ADLC Fase 0-2)}

Questo progetto è il primo dove il \texttt{\_CONTEXT.md} fa davvero la differenza tra un risultato mediocre e uno professionale.

\begin{praticabox}[title={Pratica --- Setup del progetto}]
\begin{enumerate}
  \item Crea la cartella \texttt{taskmaster}
  \item Aprila in VS~Code
  \item Crea il file \texttt{\_CONTEXT.md}:
\end{enumerate}
\end{praticabox}

\begin{lstlisting}[language={},title={\texttt{\_CONTEXT.md} per TaskMaster CLI}]
# Progetto: TaskMaster CLI

## Scopo
Applicazione CLI per la gestione di task personali con
persistenza locale. Strumento per uso quotidiano dal terminale.

## Tecnologie
- Python 3.11+
- SOLO standard library (argparse, json, os, datetime, pathlib, enum)
- Testing: pytest
- Nessuna dipendenza esterna ammessa

## Struttura del Progetto
taskmaster/
+-- _CONTEXT.md
+-- main.py              <- Entry point (solo parsing argomenti)
+-- task_manager.py      <- Classe TaskManager: logica CRUD
+-- models.py            <- Dataclass Task con tutti i campi
+-- storage.py           <- Classe Storage: lettura/scrittura JSON
+-- formatters.py        <- Formattazione output per il terminale
+-- tests/
|   +-- conftest.py      <- Fixtures condivise
|   +-- test_task_manager.py
|   +-- test_storage.py
|   +-- test_formatters.py
+-- data/
    +-- tasks.json       <- Creato automaticamente al primo uso

## Modello Dati (Task)
- id: int (auto-incrementale, non riutilizzato)
- title: str (obbligatorio, max 100 caratteri)
- description: str (opzionale, default "")
- status: enum [todo, in_progress, done] (default: todo)
- priority: enum [low, medium, high] (default: medium)
- created_at: str ISO 8601
- updated_at: str ISO 8601
- due_date: str ISO 8601 o null (opzionale)

## Comandi CLI
taskmaster add "Titolo del task"
taskmaster add "Titolo" -d "Descrizione" -p high
taskmaster add "Titolo" --due 2026-04-15
taskmaster list
taskmaster list --status todo
taskmaster list --priority high
taskmaster list --overdue
taskmaster show 3
taskmaster update 3 --status done
taskmaster delete 3
taskmaster stats

## Convenzioni di Codice
- Naming: snake_case per tutto
- Classi: PascalCase (TaskManager, Storage, Task)
- Type hints: OBBLIGATORI su tutte le firme
- Docstring: su ogni classe e metodo pubblico
- Enum: usa enum.Enum per status e priority
- Dataclass: usa @dataclass per il modello Task
- Pathlib: usa pathlib.Path, mai os.path

## Vincoli di Implementazione
- NON usare variabili globali.
- NON ignorare le eccezioni.
- NON usare exit() nel mezzo del codice. Solo in main.py.
- Il file tasks.json DEVE essere creato automaticamente.
- L'ID dei task DEVE essere auto-incrementale e MAI riutilizzato.
- Il formato JSON deve essere indentato (indent=2).
- La cancellazione di un task chiede conferma (y/N).
\end{lstlisting}

\begin{notebox}
Nota come questo \texttt{\_CONTEXT.md} è molto più dettagliato di quello del Capitolo~4. Definisce il modello dati, i comandi CLI esatti, le convenzioni di codice e i vincoli di implementazione. Più è dettagliato il contesto, migliore sarà il codice generato.
\end{notebox}

% ---------------------------------------------------------
\section{Generazione del Codice (ADLC Fase 3-4)}

\begin{praticabox}[title={Pratica --- Fase 3: Proof of Value}]
Prima di generare tutto il progetto, un test rapido. In Copilot Agent Mode:
\begin{lstlisting}[language={}]
Leggi il _CONTEXT.md. Come primo passo, genera solo il file
models.py con la dataclass Task e gli enum Status e Priority.
Non generare nient'altro per ora.
\end{lstlisting}
Verifica che usi \texttt{@dataclass}, \texttt{enum.Enum}, type hints e valori default.
\end{praticabox}

\begin{praticabox}[title={Pratica --- Fase 4: Generazione completa}]
\begin{lstlisting}[language={}]
Perfetto. Ora implementa l'intero progetto secondo il
_CONTEXT.md. Procedi in quest'ordine:
1. storage.py (persistenza JSON)
2. task_manager.py (logica CRUD)
3. formatters.py (output terminale)
4. main.py (CLI con argparse)
5. tests/ (tutti i file di test)
\end{lstlisting}
\end{praticabox}

\begin{tipbox}
Specificare l'ordine di implementazione aiuta l'IA a costruire il progetto layer per layer, evitando riferimenti circolari e dipendenze mancanti.
\end{tipbox}

% ---------------------------------------------------------
\section{Revisione del Codice Generato}

Dopo che Copilot ha generato tutti i file, fai una revisione sistematica.

\subsection*{Checklist di revisione}

\textbf{\texttt{models.py}}:
\begin{itemize}
  \item \texttt{Status} e \texttt{Priority} sono Enum, non stringhe
  \item \texttt{Task} è una dataclass con tutti i campi del contesto
  \item Metodi \texttt{to\_dict()} e \texttt{from\_dict()} per serializzazione
  \item \texttt{created\_at} e \texttt{updated\_at} generati automaticamente
\end{itemize}

\textbf{\texttt{storage.py}}:
\begin{itemize}
  \item Usa \texttt{pathlib.Path}, non \texttt{os.path}
  \item Crea il file JSON automaticamente se non esiste
  \item Legge e scrive con \texttt{indent=2}
  \item Gestisce \texttt{FileNotFoundError} e \texttt{JSONDecodeError}
\end{itemize}

\textbf{\texttt{task\_manager.py}}:
\begin{itemize}
  \item L'ID è auto-incrementale e non viene riutilizzato
  \item Supporta tutti i comandi del contesto
  \item \texttt{update\_task} modifica \texttt{updated\_at}
\end{itemize}

\textbf{\texttt{main.py}}:
\begin{itemize}
  \item Usa \texttt{argparse} con subcommand
  \item Messaggi utili, non stack trace
\end{itemize}

\textbf{\texttt{tests/}}:
\begin{itemize}
  \item Usa \texttt{tmp\_path} per lo storage
  \item Testa casi positivi e edge case
  \item Ogni test è indipendente
\end{itemize}

\begin{warnbox}
Se trovi un problema, non modificare il codice a mano. Descrivi il problema a Copilot e chiedigli di correggerlo. Questo è il flusso 0-code.
\end{warnbox}

% ---------------------------------------------------------
\section{Test ed Esecuzione}

\begin{praticabox}[title={Pratica --- Test}]
\begin{lstlisting}[language=bash]
python -m pytest tests/ -v
\end{lstlisting}
Se qualche test fallisce, copia l'errore nella chat di Copilot:
\begin{lstlisting}[language={}]
Questo test fallisce con il seguente errore:
[incolla l'errore]
Correggi il codice per far passare il test.
\end{lstlisting}
\end{praticabox}

\begin{praticabox}[title={Pratica --- Uso dell'applicazione}]
\begin{lstlisting}[language=bash]
# Aggiungi task
python main.py add "Leggere il capitolo 6" -p high
python main.py add "Comprare il latte" --due 2026-04-01

# Lista tutti
python main.py list

# Filtra
python main.py list --status todo
python main.py list --priority high

# Aggiorna e completa
python main.py update 1 --status in_progress
python main.py update 1 --status done

# Statistiche
python main.py stats
\end{lstlisting}
\end{praticabox}

\begin{checkpointbox}
Se tutti i comandi funzionano e i test passano, hai un'applicazione CLI professionale completa.
\end{checkpointbox}

% ---------------------------------------------------------
\section{Iterazione: Aggiungere Funzionalità}

\begin{praticabox}[title={Pratica --- Esporta in formato Markdown}]
\begin{lstlisting}[language={}]
Aggiungi un comando "export" che esporta tutti i task in un
file Markdown formattato come tabella.
Uso: python main.py export report.md
Il file deve contenere una tabella con i task organizzati
per priorita' (high -> medium -> low).
Aggiungi i test corrispondenti.
\end{lstlisting}
\end{praticabox}

\begin{praticabox}[title={Pratica --- Ricerca per testo}]
\begin{lstlisting}[language={}]
Aggiungi un comando "search" che cerca tra titoli e
descrizioni. Uso: python main.py search "email"
La ricerca deve essere case-insensitive.
Aggiungi i test corrispondenti.
\end{lstlisting}
\end{praticabox}

% ---------------------------------------------------------
\section{Il Pattern che Userai Sempre}

Questo progetto ha stabilito il pattern che seguirai per tutto il libro:

\begin{enumerate}
  \item \textbf{\texttt{\_CONTEXT.md}} $\rightarrow$ Definisci il progetto (5--15~min)
  \item \textbf{Proof of Value} $\rightarrow$ Genera un pezzo piccolo e verifica (5~min)
  \item \textbf{Generazione} $\rightarrow$ L'IA implementa tutto (10--20~min)
  \item \textbf{Revisione} $\rightarrow$ Controlli la qualità (5--10~min)
  \item \textbf{Test} $\rightarrow$ Verifica che funzioni (2--5~min)
  \item \textbf{Iterazione} $\rightarrow$ Aggiungi funzionalità e correggi (continuo)
  \item \textbf{Aggiornamento} $\rightarrow$ Migliora il \texttt{\_CONTEXT.md} (2~min)
\end{enumerate}

Questo pattern scala: funziona per un CLI tool da 200 righe come per un'app full-stack da 10.000.

% ---------------------------------------------------------
\section{Lezioni Apprese (ADLC Fase 6)}

\begin{lstlisting}[language={},title={Lezioni da aggiungere al contesto}]
## Lezioni Apprese
- Specificare l'ordine di implementazione
  (storage -> logic -> UI -> main -> tests)
- I test con tmp_path devono creare un'istanza fresh
- argparse con subparsers richiede set_defaults(func=...)
- L'IA tende a dimenticare la conferma y/N sulla
  cancellazione: metterla nei vincoli
\end{lstlisting}

% ---------------------------------------------------------
\section*{Riepilogo del Progetto}

\begin{center}
\begin{tabularx}{\textwidth}{l X}
\toprule
\textbf{Aspetto} & \textbf{Dettaglio} \\
\midrule
Linguaggio & Python 3.11+ \\
Dipendenze & Zero (solo standard library) \\
File generati & $\sim$7 file, $\sim$400--600 righe totali \\
Test & $\sim$15--20 test case \\
Tempo totale & $\sim$30--45 minuti \\
Codice scritto a mano & Solo il \texttt{\_CONTEXT.md} \\
\bottomrule
\end{tabularx}
\end{center}
